\documentclass{article}
\usepackage[UTF8]{ctex}
\usepackage{amsmath,amssymb,amsfonts}
\usepackage{enumitem}
\usepackage{geometry}
\geometry{a4paper, margin=1in}
% === 错题本排版优化（不参与题目标识） ===
\setlength{\parindent}{0pt}
\setlength{\parskip}{0.35em plus 0.10em minus 0.08em}
\setlength{\emergencystretch}{3em}
\setlist{itemsep=0.25em, topsep=0.40em, parsep=0pt}
\allowdisplaybreaks[2]
% === 错题本排版优化结束 ===
\begin{document}
\section*{电子科技大学2024年高等代数真题}
\section*{158．电子科技大学2024年高等代数试卷（回忆版）}
一、填空题。\\

\subsection{题目1}
1．已知矩阵 \(A=\left(\begin{array}{cc}1 & -1 \\ 1 & 1\end{array}\right), B=\frac{\sqrt{3}}{2} A\) ，则 \(\left(\begin{array}{cc}A & O \\ O & B\end{array}\right)^{\mathrm{n}}=\) \(\_\_\_\_\)．\\
\(\_\_\_\_\)

\subsection{题目2}
2．已知 \(G=\left(\eta_{1}, \eta_{2}, \eta_{3}\right)\) 为正交矩阵，则矩阵 \(A=2 \eta_{1} \eta_{1}^{\mathrm{T}}+\eta_{2} \eta_{2}^{\mathrm{T}}+\eta_{3} \eta_{3}^{\mathrm{T}}\) 的最小多项式为 \(\_\_\_\_\)．

\subsection{题目3}
3．已知矩阵 \(A\) 满足 \(A^{4}=O\) ，则 \(A\) 按照相抵分类一共可分为 \(\_\_\_\_\)类．

\subsection{题目4}
4．已知矩阵 \(A\) 为 \(m \times n\) 矩阵，且 \(r(A)=r<n\) 。则 \(A X=b\) 的解中线性无关的向量个数最多为\\
\(\_\_\_\_\)个．

\subsection{题目5}
5．设 \(A\) 为4阶实对称矩阵，且三个特征值为 \(\lambda_{1}=1, \lambda_{2}=3, \lambda_{3}=4\) ，且 \(|A|=-12\) ，其中 \(V_{1}, V_{2}, V_{3}\)分别为 \(\lambda_{1}, \lambda_{2}, \lambda_{3}\) 的特征子空间，则 \(\operatorname{dim}\left(\left(V_{1} \oplus V_{2} \oplus V_{3}\right)^{\perp}\right)=\) \(\_\_\_\_\)。

\subsection{题目6}
6．已知 \(\Phi\) 是 \(\operatorname{End}_{F}(V)\) 上的线性变换，并且多项式 \(f(x)\) 无重根，则 \(\operatorname{Im} f(\Phi)+\operatorname{Im} f^{\prime}(\Phi)=\) \(\_\_\_\_\)．
\end{document}
